
\documentclass[authoryear,preprint,review,times,10pt]{elsarticle}

\usepackage{amsmath}    
\usepackage{amsthm}    
\usepackage{amssymb}     
\usepackage{amsfonts}  
\usepackage{bm}         

\usepackage{booktabs}   
\usepackage{multirow}    
\usepackage{longtable}   
\usepackage{threeparttable} 
\usepackage{tabularx}    
\usepackage{caption}      
\usepackage{adjustbox}    
\usepackage{pdflscape}    
\usepackage{enumitem}     
\usepackage{lineno}       
\usepackage{setspace}     
\doublespacing           

\usepackage{graphicx}
\usepackage{subfig}
\graphicspath{{simulation/figures/}}
\usepackage{tikz}
\usetikzlibrary{positioning, shapes.geometric, arrows}

\usepackage{geometry}
\geometry{
    left=2.5cm,
    right=2.5cm,
    top=2.5cm,
    bottom=2.5cm,
}

\usepackage{hyperref}    

\makeatletter
\def\ps@pprintTitle{
  \let\@oddhead\@empty
  \let\@evenhead\@empty
  \def\@oddfoot{\hfil\thepage\hfil}
  \let\@evenfoot\@oddfoot}
\def\pprintMaketitle{\clearpage
  \iflongmktitle\if@twocolumn\let\columnwidth=\textwidth\fi\fi
  \resetTitleCounters
  \def\baselinestretch{1}
  \printFirstPageNotes 
  \begin{center}
   \thispagestyle{pprintTitle}
   \def\baselinestretch{1}
    \Large\@title\par\vskip18pt
    \normalsize\elsauthors\par\vskip10pt
    \normalsize\itshape\elsaddress\par\vskip12pt
    \ifvoid\absbox\else\unvbox\absbox\par\vskip10pt\fi
    \ifvoid\keybox\else\unvbox\keybox\par\vskip10pt\fi
  \end{center}
  \gdef\thefootnote{\arabic{footnote}}%
  }
\makeatother

\journal{}

\begin{document}

\newtheorem{prop}{Proposition}

\modulolinenumbers[1]
\runninglinenumbers
\begin{frontmatter}

\title{Cooperation in competition? A biform game model for waste concrete recycling considering the cap-and-trade mechanism}

\author[1]{Zhe Zhang}
\author[2]{Qinghua He, Ph.D.}
\author[3]{Zilun Wang, Ph.D.}

\affiliation[1]{Postgraduate Student, School of Economics and Management, Tongji University, Shanghai 200092, China; Email: 2431220@tongji.edu.cn}
\affiliation[2]{Professor, School of Economics and Management, Tongji University, Shanghai 200092, China; Email: heqinghua@tongji.edu.cn}
\affiliation[3]{Lecturer, School of Civil Engineering, Suzhou University of Science and Technology, Suzhou 215011, China (Corresponding author); Email: zylenw97@usts.edu.cn}


\begin{abstract}

  Waste concrete recycling is essential for sustainable construction and resource management. Concrete manufacturers and recyclers face multiple challenges as the participants. They engage in pricing competition and have the option to pursue research and development investment cooperation. However, existing research typically treats competition and cooperation as isolated stages, leaving it unclear whether the enterprises should cooperate within a competitive environment and what the resulting economic benefits might be. Additionally, the cap-and-trade mechanism acts as a crucial policy tool for balancing environmental protection and economic development. This complicates their decisions and financial interests. To address these gaps, we develop a biform game model for waste concrete recycling that simultaneously integrates cooperative and non-cooperative games under the influence of the cap-and-trade mechanism. We compare this with a pure non-cooperative game model. We find that cooperation incentivizes a higher conversion rate of the waste concrete and leads to an expansion of total collection quantities. Cooperation in the competition environment achieves a Pareto improvement, where both the manufacturer and the recycler realize higher individual profits compared to the non-cooperative game. We also find that government financial incentives are not universally beneficial. Excessive subsidies can destabilize the cooperative alliance. Furthermore, the cap-and-trade mechanism raises costs and reduces profit margins for both participants, but cooperation provides a buffering effect against such volatility. This study expands the application of biform games and contributes to the literature on waste concrete recycling management and economic decision making. It also provides essential managerial insights for enterprises seeking cooperation in the competitive environment and for government policy design in the carbon market.

\end{abstract}

\begin{keyword}
Waste concrete recycling \sep Biform Game \sep Cap-and-trade Mechanism \sep Cooperation
\end{keyword}

\end{frontmatter}

\section*{Managerial relevance statement}
This study develops a biform game model to investigate manufacturer-recycler co-opetition under the cap-and-trade mechanism, yielding actionable guidance for practitioners and policymakers through equilibrium and numerical analyses. First, concrete manufacturers and recyclers should actively seek R\&D cooperation, as pure price competition in the collection market fails to maximize mutual economic interests. Second, enterprise managers should evaluate the operational impacts of environmental policies like the cap-and-trade mechanism. To effectively buffer against carbon market volatility and high compliance costs, firms are advised to implement KPI-linked technology cost-sharing contracts and establish internal carbon risk-sharing mechanisms, transforming regulatory constraints into cooperative innovation drivers. Finally, government policymakers should simultaneously weigh enterprise economic viability and long-term environmental development when setting carbon emission caps and trading prices. Furthermore, rather than simply increasing financial incentives, the government should carefully evaluate and maintain subsidies within an optimal range, using targeted structural support like logistical subsidies to stimulate the market while fostering a sustainable R\&D cooperative ecosystem between concrete manufacturers and recyclers.


\section{Introduction}
The recycling and processing of renewable resources has become one challenge in engineering management, with the construction industry serving as one of the most significant sectors for these activities \citep{yuan2025Exploring}. Rapid development of the construction industry in recent decades has led to a sharp increase in construction and demolition (C.and.D) waste generation, with waste concrete accounting for more than half of total C.and.D waste \citep{zhao2020Use}. Given the environmental and social damage caused by this waste, urgent management of its recycling is required \citep{wu2019Offsite}. Waste concrete recycling has demonstrated potential to alleviate environmental pressures and shortages of construction raw materials due to its regenerative and utilization value \citep{mostert2021Climate}. Many countries, particularly developed nations with limited natural resources, have widely adopted recycled concrete made by the waste concrete \citep{yenipazarli2019Incentives}. However, promotion of waste concrete recycling management faces obstacles including immature technologies, unclear economic benefits, and incomplete policies, resulting in low enterprise participation \citep{liu2021Explore,ma2020Challenges}. The waste concrete recycling process involves two major participants: concrete manufacturers and recyclers \citep{he2020Investigation,lu2015Stakeholders}. Confronted with high technical costs and market uncertainties, concrete manufacturers and recyclers often lack the capacity to independently overcome these barriers and require inter-enterprise coordination to achieve complementary advantages, risk sharing, and maximum economic benefits from waste concrete recycling \citep{bao2019Procurement}. Therefore, how to effectively coordinate the waste concrete recycling management process has become a critical management challenge constraining the development of the waste concrete recycling industry.

The waste concrete recycling process generally involves collecting raw waste concrete from the market and the treatment \citep{he2020Investigation}. In this process, concrete manufacturers possess the channel advantage. They can collect waste concrete from construction sites using their delivery trucks on return trips, which keeps their transportation costs relatively low \citep{besklubova2026Establishing}. On the other hand, the recyclers use dedicated equipment to process the collected waste into recycled aggregates\citep{chen2025Exploring}. During these interactions, the two participants engage in price competition at the front end to secure limited waste concrete resources \citep{he2020Investigation}. At the same time, they can choose to cooperate at the back end by sharing the research and development investment required to improve these processing technologies \citep{liu2019Expand}. These intertwined behaviors form a complex co-opetition relationship. However, both participants are often hesitant to actually commit to the R\&D investment cooperation \citep{bao2020Developing}. This is because their competitive and cooperative behaviors mutually influence each other. Pricing battles at the front end impact the final profit distribution from R\&D investment sharing at the back end, making the ultimate economic benefits of working together unclear \citep{li2020Key}. While previous research has explored recycling strategies, most studies treat collection and treatment or competition and cooperation as separate stages \citep{hu2021Dynamic,yuan2023Differentiated}. They fail to clarify whether these companies should cooperate in such a competitive environment and how they should adjust their prices based on this cooperation. Therefore, it is necessary to explore how the interaction between pricing competition and R\&D investment cooperation affects the decisions and economic benefits for both concrete manufacturers and recyclers.

The environmental impact of the recycling process also complicates the co-opetition relationships. Although waste concrete recycling offers advantages in conserving virgin resources, the treatment process is often accompanied by substantial carbon emissions \citep{difilippo2019Impacts}. Current mainstream waste concrete recycling technologies primarily rely on crushing, separation, and calcination-based thermal treatment, which still generate considerable carbon dioxide \citep{oh2021Proposal}. To curb industrial carbon emissions, the cap-and-trade mechanism has been widely implemented globally as an effective market-based tool that balances economic benefits with environmental protection \citep{oh2021Proposal,yenipazarli2019Incentives}. Waste concrete recycling is also an industry widely affected by the cap-and-trade mechanism.
For the waste concrete recycling industry, carbon quota constraints and fluctuations in trading prices reshape the cost structures and profit margins of participating enterprises, compelling them to consider the economic implications of carbon emissions when making decisions \citep{zheng2025Coopetition}. However, existing research on waste concrete recycling management and economic decision-making has largely overlooked this aspect. The impacts of the cap-and-trade mechanism on these co-opetition relationships remain unclear, resulting in a lack of clarity on how concrete manufacturers and recyclers should adjust their decisions under such environmental policy constraints. Therefore, incorporating the cap-and-trade mechanism into the research framework to investigate how it influences the economic benefits of waste concrete recycling and the decision-making behavior of enterprises is necessary.

This study proposes a two-stage biform game model for waste concrete recycling that incorporates both non-cooperative and cooperative strategies to analyze the co-opetition relationships of concrete manufacturers and recyclers. Non-cooperative strategies represent recycling pricing competition between waste manufacturers and recyclers, which do not yield final equilibrium solutions but rather establish a competitive context for their R\&D investment cooperation decisions. Simultaneously, cooperative strategies are jointly determined by both enterprises. The profit distribution resulting from the cooperative strategies subsequently becomes the critical profit function for resolving the final non-cooperative strategies. Furthermore, we compare this co-opetition strategy with a non-cooperative game model. Specifically, this study addresses the following research questions:

\begin{itemize}
\item \textbf{RQ1:} Considering cap-and-trade mechanism, should concrete manufacturers and recyclers engage in cooperation within the competitive environment?
\item \textbf{RQ2:} What are the impacts of the cap-and-trade mechanism on the decisions of concrete manufacturers and recyclers in waste concrete recycling?
\end{itemize}

The remainder of this paper is organized as follows. Section 2 reviews the relevant literature. Section 3 presents the model formulation, including problem description, assumptions, and profit functions. Section 4 derives the equilibrium analysis under non-cooperative and cooperative game scenarios. Section 5 conducts numerical simulations to verify the theoretical results and analyze the effects of key parameters. Section 6 discusses the findings and implications. Section 7 concludes the paper with recommendations for future research.

\section{Literature review}
\subsection{Waste concrete recycling management and economic decision}
In the waste concrete recycling supply chain, existing literature exploring their management and economic decision can be categorized into three main research themes: reverse logistics network design and operation management, government subsidy and incentive mechanism design, and manufacturer-recycler competition and cooperation relationships.

Regarding reverse logistics network design and operation management, researchers primarily focus on optimizing physical flows and capacity. \cite{rahimi2018Sustainable} developed a multi-period mixed integer linear programming (MILP) model to optimize the network design for C.and.D waste recycling. Similarly, \cite{xu2019Reverse} proposed a capacity expansion pricing mechanism tailored for recycling networks under demand uncertainty , while \cite{liu2019Expand} identified key cost drivers through a comprehensive network cost analysis. However, these operational models treat the recycling network as a centrally coordinated system, overlooking the decentralized decision-making processes and the profit-driven strategic conflicts between independent concrete manufacturers and recyclers.

Within the theme of government subsidy and incentive mechanism design, studies emphasize utilizing external policy interventions to stimulate market participation. \cite{yuan2023Differentiated} investigated differentiated subsidy schemes under heterogeneous consumer quality perceptions, demonstrating that uniform subsidies often fail to achieve optimal outcomes. \cite{su2024Stakeholder} evaluated the impact of government intervention through quota trading mechanisms on recycling rates , and \cite{zheng2024Optimal} examined how various incentive combinations influence stakeholder participation. 
Although existing research confirms the effectiveness of macro-level policy tools, it often assumes uniform cost structures among supply chain participants. The literature inadequately explains how enterprises' specific operational advantages dictate varying responses to environmental policies.

Finally, concerning manufacturer-recycler the competition and cooperation relationships of the concrete manufacturer and the recycler, scholarship has begun to unpack micro-level strategic behaviors. \cite{he2020Investigation} analyzed optimal decisions under different market scenarios to identify the conditions under which cooperation dominates pure competition. \cite{lin2025Cooperation} revealed the cooperation dilemma (prisoner's dilemma) in C.and.D waste recycling, showing that unilateral strategies frequently lead to suboptimal industry outcomes. Furthermore, \cite{li2025Mechanism} employed evolutionary game theory to study the long-term dynamics of manufacturer-recycler relationships. While these contributions provide essential theoretical guidance for the decision-making of concrete manufacturers and recyclers in waste concrete recycling, existing research predominantly focuses on single-scenario analyses or separate scenario comparisons. They fail to capture the simultaneous interaction patterns of competition and cooperation inherent in this process, struggling to resolve the hesitation of R\&D investment cooperation within the industry.

\subsection{Biform games and recycling supply chain coordination}
The theoretical foundation of this study rests on the biform game model, which was originally formalized by \cite{brandenburger2007Biform} to analyze scenarios where firms simultaneously engage in competition and cooperation. This model captures the coopetition relationships by integrating non-cooperative game theory (first stage) with cooperative game theory (second stage). In the first stage, enterprises make strategic moves to shape the competitive environment; in the second stage, they allocate the created surplus through bargaining mechanisms, typically using mechanisms such as the Shapley value.

The application of biform games to supply chain coordination has received substantial attention in recent years. \cite{feess2014Surplus} developed a pioneering supply chain biform game model where firms make non-cooperative investment decisions in the first stage and divide surplus via the Shapley value in the second stage, revealing that all firms' investment incentives fall below the socially optimal level. \cite{fiala2022Modelling} applied biform games to model co-opetition in network industries, distinguishing between sequential and simultaneous game structures. More recently, \cite{zheng2023New} proposed a biform game-based investment incentive mechanism for eco-efficient innovation in supply chains.

Extending this framework to reverse supply chain and recycling contexts, several studies have examined how enterprises coordinate collection and treatment activities under competitive pressures. \cite{jia2023Novel} employed a biform game to coordinate a multi-agent reverse supply chain with fairness concerns, demonstrating that the biform game mechanism could improve all supply chain members' outcomes compared to pure non-cooperative equilibria. \cite{gui2018Design} analyzed design incentives under collective extended producer responsibility using biform games, focusing on how processor capacity sharing affects design decisions and alliance stability in recycling networks. \cite{wang2026Integration} explicitly modeled collection competition and processing cooperation by constructing five game scenarios in a reverse supply chain with horizontal competitive recycling, revealing that full cooperation maximizes total system profit while different allocation methods favor different participants. The integration of biform games with environmental policy mechanisms has also been explored. \cite{zheng2024New} developed a biform game-based coordination mechanism for a carbon complementary supply chain under hybrid carbon regulations, identifying threshold conditions for regulatory intensity beyond which cooperative investment yields Pareto improvements.

Therefore, the biform game framework is well-suited for analyzing waste concrete recycling, where competition in collection markets coexists with cooperation in treatment activities. This analytical approach enables simultaneous examination of how carbon cap-and-trade mechanisms influence collection pricing strategies and how R\&D investment cooperation can be designed to improve individual profits for both concrete manufacturers and recyclers. We develop a biform game model tailored to the waste concrete recycling context to explore their optimal economic decisions.

\subsection{The cap-and-trade mechanism and recycling}
The cap-and-trade mechanism sets a carbon emission cap while allowing enterprises to freely trade emission allowances, incentivizing enterprises to pursue lower-cost emission reduction pathways while providing a flexible balance between economic benefits and environmental protection \citep{li2019How}.The implementation of cap-and-trade mechanisms has reshaped the cost structures and operational boundaries of recycling supply chains. A primary stream of literature evaluates the effectiveness of this market-based instrument in balancing economic profitability with environmental compliance. Studies have demonstrated that cap-and-trade policies, when appropriately parameterized, can effectively incentivize emission reductions without severely compromising system performance \citep{bai2022Distributionally}. Furthermore, comparative analyses indicate that the optimal regulatory choice—whether carbon taxes, strict caps, or cap-and-trade—is highly contingent upon market structure, carbon price levels, and firms' technological innovation capacities \citep{lyu2022Manufacturers,huang2024Optimal}. These foundational studies confirm that carbon quotas are no longer mere exogenous constraints, but endogenized variables that dictate recycling participation.

Building on this macro-policy premise, recent scholarship has increasingly focused on how carbon trading influences strategic interactions and mode selection in recycling networks. Under the dual pressures of economic profit and carbon compliance, researchers have highlighted the necessity of joint recycling modes. For instance, dynamic carbon pricing has been shown to induce symbiotic cooperation between manufacturers and recyclers, often outperforming solitary recycling efforts under policy interventions \citep{cao2023Developing,wu2024Electric}. Additionally, optimal recycling channel selection and supply chain pricing strategies are heavily modulated by the allocation of carbon quotas \citep{zhang2022Recycling}.

Despite these valuable insights, a misalignment exists between current carbon-constrained recycling models and the engineering realities of waste concrete. Existing literature predominantly investigates high-value-density, low-transport-friction products (e.g., electronic waste and EV batteries). In contrast, waste concrete is characterized by low unit value, massive physical volume, and highly asymmetrical transportation costs (i.e., the concrete manufacturer's backhaul advantage versus the recycler's dedicated logistics). Furthermore, the processing of waste concrete is inherently carbon-intensive, making the recycler highly sensitive to carbon price volatility. Current models fail to capture how the cap-and-trade mechanism specifically couples with these unique logistical asymmetries to influence front-end pricing competition and back-end R\&D investment. Consequently, extending the cap-and-trade analysis to the waste concrete context recycling is imperative to formulate effective, industry-specific economic strategies.

\section{Model formulation}
\subsection{Problem description}
We consider a two-stage waste concrete recycling supply chain consisting of a traditional concrete manufacturer M and a specialized waste concrete recycler R. Both M and R participate in the collection of waste concrete from waste generators, while recycler R possesses specialized technologies and equipment to process the collected waste into recycled concrete aggregate (RCA). Recycler R generates carbon dioxide emissions during this operations, and therefore must consider the impact of cap-and-trade policies on its economic decisions. In this section, we develop and compare two different game-theoretic models: (1) a pure non-cooperative waste concrete recycling model considering the cap-and-trade mechanism, and (2) a biform game model combining R\&D investment cooperation with the cap-and-trade mechanism.

In the non-cooperative game model, there is no R\&D investment sharing between the concrete manufacturer and the recycler(as shown in Fig. 1(a)). Recycler R independently bears all the R\&D investment for upgrading the processing technology. The decision sequence is as follows: first, recycler R independently determines the optimal RCA conversion rate to maximize its own profit; subsequently, based on this established conversion rate, manufacturer M and recycler R simultaneously determine their respective collection prices to compete for waste concrete resources in the market.

In the biform game model, the interaction involves both pricing competition and R\&D investment cooperation(as shown in Fig. 1(b)). The game process is modeled in two stages. First, in the non-cooperative stage, concrete manufacturer M and recycler R determine their respective collection prices to form a competitive situation. Second, in the cooperative stage, under any given competitive situation, the two participants jointly form an alliance to determine the optimal conversion rate and the R\&D investment sharing proportion. We calculate the characteristic values of each coalition and apply the Shapley value method to determine the profit distribution values ($\varphi_M$ and $\varphi_R$) for manufacturer M and recycler R, respectively. Finally, using these profit distribution values as the payoff functions, we solve the first-stage non-cooperative game backward to derive the optimal collection prices, conversion rate, and final profit allocations.

The following assumptions are made:

\textbf{Assumption 1.} Manufacturer $M$ and recycler $R$ set collection prices $(p_M, p_R)$ for waste concrete generators. Since both parties participate in collection, there exists competitive pressure in the recycling market. Let $\beta$ denote the competition intensity between $M$ and $R$. Additionally, the collection quantity depends on the collection price, with $\alpha$ representing the sensitivity of waste generators to collection prices. We assume sufficient waste concrete available in the market and $R$ has sufficient capacity to process all collected quantity. Referring to \citep{chen2025Lowcarbon} and \citep{zhang2024Selection}'s research,the collection quantity function follows:
\begin{equation}
q_i = \alpha p_i - \beta p_j \quad (i,j \in \{M, R\}, i \neq j)
\end{equation}
where $\alpha > \beta > 0$ ensures positive collection quantities. A larger $\alpha$ indicates higher sensitivity to prices and thus greater collection quantity, and a larger $\beta$ indicates more intense competition between $M$ and $R$.

\textbf{Assumption 2.} The recycling processing capability of recycler $R$ is characterized by the RCA conversion rate $k \in [0,1]$. Improving this conversion rate (e.g., through enhanced mortar removal technology or multi-stage air separation) requires fixed investment in research and development. Based on the research by \cite{jones2011Information} and \cite{chen2024Risk}, we assume that the R\&D investment cost function is:
\begin{equation}
C(k) = \frac{1}{2} c k^2
\end{equation}
where $c$ is the R\&D investment cost coefficient.

\textbf{Assumption 3.} In practice, traditional concrete manufacturers often lack deep-processing environmental equipment or are constrained by site and environmental qualifications. Therefore, manufacturer $M$ outsources the collected waste concrete to recycler $R$ for centralized processing. Manufacturer $M$ pays a unit outsourcing fee $w$ to recycler $R$, while recycler $R$ incurs actual operational costs $m$ (e.g., utilities, equipment depreciation, labor) per unit of waste concrete processed.

\textbf{Assumption 4.} Assume $\theta$ is the market price per unit of recycled concrete aggregate (RCA) when perfectly converted from waste concrete. Under the actual RCA conversion rate $k$ achieved by recycler $R$, the actual extracted value per unit of waste concrete is $\theta k$. For manufacturer $M$, this value represents the cost saving from substituting natural aggregates with RCA. For recycler $R$, it represents the selling price of processed RCA. To encourage construction waste recycling under the zero-waste city initiative, the government provides a subsidy $s$ per unit of waste concrete collected, which is distributed to the respective collector based on their collection quantity.

\textbf{Assumption 5.} To obtain high-quality recycled aggregate and maintain stable outsourcing cooperation, manufacturer $M$ is willing to share the R\&D investment for upgrading the processing technology with recycler $R$. Let manufacturer $M$ share a proportion $(1-n)$ and recycler $R$ bear a proportion $n$, where $n \in [0,1]$. When $n=1$, all investment costs are borne solely by recycler $R$.

\textbf{Assumption 6.} Both manufacturer $M$ and recycler $R$ incur unit transportation costs for collecting waste concrete from generators\citep{chileshe2016Analysis}. The physical properties of waste concrete (heavy, large volume, low value density) make transportation costs a significant component of total recycling costs. Due to the backhaul logistics advantage in the construction industry, manufacturer $M$'s concrete delivery trucks can collect waste concrete on return trips from construction sites, resulting in a lower transportation cost $\mu_M$. In contrast, recycler $R$ must dispatch dedicated vehicles for collection, incurring a higher transportation cost $\mu_R$. We assume $\mu_M < \mu_R$. This cost asymmetry represents the engineering management characteristic of waste concrete recycling.

\textbf{Assumption 7.} Under the cap-and-trade mechanism, recycler $R$ is subject to government-allocated carbon emission quotas and engages in carbon trading. Let $G$ denote the total carbon emission quota allocated by the government to recycler $R$ in a compliance cycle. The actual carbon emissions generated by recycler $R$ processing a unit of waste concrete is denoted by $e$ (treated as a constant). The total processing quantity is $q_{total} = q_M + q_R = (\alpha - \beta)(p_M + p_R)$. Let $p_c$ denote the trading price per unit of carbon emission in the carbon market. The net carbon trading cost incurred by recycler $R$ is:
\begin{equation}
p_c \left[ e \cdot q_{total} - G \right] = p_c \left[ e(\alpha - \beta)(p_M + p_R) - G \right]
\end{equation}

\begin{table}[!htbp]
\centering
\captionsetup{font=normalsize, labelsep=period}
\caption{Notation and definitions}
\label{tab:notation}
\normalsize
\begin{threeparttable}
\begin{tabular*}{0.9\textwidth}{@{\extracolsep{\fill}}lc}
\toprule
\textbf{Symbol} & \textbf{Definition} \\
\midrule
$\theta$ & Unit potential value of waste concrete converted to highest-quality RCA \\
$\alpha$ & Sensitivity of waste generators to collection price \\
$\beta$ & Competition intensity between manufacturer and recycler, $\beta \in [0,1]$ \\
$\mu_M$ & Unit transportation cost of manufacturer $M$ \\
$\mu_R$ & Unit transportation cost of recycler $R$ \\
$p_M$ & Unit collection price paid by manufacturer $M$ to waste generators \\
$p_R$ & Unit collection price paid by recycler $R$ to waste generators \\
$q_M$ & Waste concrete collection quantity of manufacturer $M$ \\
$q_R$ & Waste concrete collection quantity of recycler $R$ \\
$n$ & Proportion of R\&D investment cost borne by recycler $R$, $n \in [0,1]$ \\
$k$ & RCA conversion rate of recycler $R$, $k \in [0,1]$ \\
$c$ & R\&D investment cost coefficient \\
$m$ & Actual operational cost of recycler $R$ per unit waste concrete \\
$s$ & Government subsidy per unit waste concrete processed \\
$w$ & Unit outsourcing fee paid by manufacturer $M$ to recycler $R$ \\
$G$ & Total carbon emission quota allocated by government to recycler $R$ \\
$e$ & Actual carbon emissions per unit of waste concrete processed by recycler $R$ \\
$p_c$ & Trading price per unit of carbon emission in the carbon market \\
$\Pi_i$ & Profit function of participant $i$, $i \in \{M,R\}$ \\
$\varphi_i$ & Shapley value (profit allocation) for participant $i$, $i \in \{M,R\}$ \\
\bottomrule
\end{tabular*}
\begin{tablenotes}[flushleft]
\small\linespread{1}\selectfont
\item \textit{Note}: All decision variables and parameters are non-negative unless otherwise specified.
\end{tablenotes}
\end{threeparttable}
\end{table}
\vspace{10pt}
\normalsize

Based on the above assumptions, the profit functions are expressed as:
\begin{equation}
\Pi_M = (\theta k - p_M - w - \mu_M + s)(\alpha p_M - \beta p_R) - (1 - n)\frac{c k^2}{2}
\end{equation}
\begin{equation}
\Pi_R = (\theta k - p_R - m - \mu_R + s)(\alpha p_R - \beta p_M) + (w - m)(\alpha p_M - \beta p_R) - n\frac{c k^2}{2} - p_c\left[ e(\alpha - \beta)(p_M + p_R) - G \right]
\end{equation}

\subsection{Non-cooperative game module}
In this module, there is no R\&D investment sharing between the manufacturer and the recycler. Specifically, the recycler $R$ bears all the R\&D investment costs ($n=1$), while the manufacturer $M$ does not participate in the investment. Under this setting, the recycler independently optimizes the conversion rate $k$ to maximize its own profit, and subsequently both participants compete in collection prices.

Under $n=1$, the profit functions simplify to:
\begin{equation}
\Pi_M^N = (\theta k - p_M - w - \mu_M + s)(\alpha p_M - \beta p_R)
\label{eq:profit_MN}
\end{equation}
\begin{equation}
\Pi_R^N =
\begin{aligned}
& (\theta k - p_R - m - \mu_R + s)(\alpha p_R - \beta p_M) + (w - m)(\alpha p_M - \beta p_R) \\
& - \frac{c k^2}{2} - p_c\left[ e(\alpha - \beta)(p_M + p_R) - G \right]
\end{aligned}
\label{eq:profit_RN}
\end{equation}

We solve this game using backward induction. The proof of Proposition 1 is provided in Appendix A. We now present the main result as follows.

\begin{prop}
\label{prop:baseline}
Under the non-cooperative game module (baseline), the equilibrium strategies are given by:
\begin{enumerate}[label={(\arabic*)}]
\item The optimal conversion rate chosen by the recycler is:
\begin{equation}
k^{N*} = \frac{\theta(\alpha p_R^{N*} - \beta p_M^{N*})}{c}
\label{eq:k_N_star}
\end{equation}

\item The equilibrium collection prices $(p_M^{N*}, p_R^{N*})$ satisfy:
\begin{equation}
p_M^{N*} = \frac{
\displaystyle
\begin{aligned}
& 4c^2 \left[ 2\alpha^2(s - w - \mu_M) + \alpha\beta(s - m - \mu_R) - \beta^2(w - m) \right] \\
& - 2c\theta^2 \Big\{ (\alpha-\beta) \left[ \alpha^2(s - w - \mu_M) - (\alpha^2 - \alpha\beta + \beta^2)(w - m) \right] - \alpha(\alpha^2 - \alpha\beta + \beta^2)(\mu_M - \mu_R) \Big\} \\
& - 4c\alpha p_c e (\alpha - \beta)(\alpha^2 - \alpha\beta + \beta^2)
\end{aligned}
}{
\displaystyle (\alpha-\beta)^3(\alpha+\beta)\theta^4 - 4c(\alpha-\beta)(3\alpha^2-\beta^2)\theta^2 + 4c^2(4\alpha^2-\beta^2)
}
\label{eq:p_MN_star}
\end{equation}

\begin{equation}
p_R^{N*} = \frac{
\displaystyle
\begin{aligned}
& 4c^2\alpha \left[ 2\alpha(s - m - \mu_R) + \beta(s - 3w + 2m - \mu_M) \right] \\
& - 2c\alpha\theta^2 \Big\{ (\alpha-\beta) \left[ \alpha(w - m) + \beta(s - 3w + 2m - \mu_M) \right] + \alpha(\alpha - 2\beta)(\mu_M - \mu_R) \Big\} \\
& - 4c p_c e (\alpha - \beta) \left[ c(2\alpha^2 - \beta^2) - \theta^2\alpha(\alpha - \beta) \right]
\end{aligned}
}{
\displaystyle (\alpha-\beta)^3(\alpha+\beta)\theta^4 - 4c(\alpha-\beta)(3\alpha^2-\beta^2)\theta^2 + 4c^2(4\alpha^2-\beta^2)
}
\label{eq:p_RN_star}
\end{equation}

\item The equilibrium collection quantities are:
\begin{equation}
q_M^{N*} = \alpha p_M^{N*} - \beta p_R^{N*}
\quad \text{and} \quad
q_R^{N*} = \alpha p_R^{N*} - \beta p_M^{N*}
\label{eq:q_N_star}
\end{equation}

\item The equilibrium profits are:
\begin{equation}
\Pi_M^{N*} = (\theta k^{N*} - p_M^{N*} - w - \mu_M + s)q_M^{N*}
\label{eq:Pi_MN_star}
\end{equation}
\begin{equation}
\Pi_R^{N*} = (\theta k^{N*} - p_R^{N*} - m - \mu_R + s)q_R^{N*} + (w - m)q_M^{N*}
+ \frac{\theta^2(q_R^{N*})^2}{2c} - p_c\left[ e(\alpha - \beta)(p_M^{N*} + p_R^{N*}) - G \right]
\label{eq:Pi_RN_star}
\end{equation}
\end{enumerate}
\end{prop}

 The proof is provided in Appendix A.

Note that the non-cooperative game model characterizes the case where there is no R\&D investment sharing between the concrete manufacturer and the recycler. The recycler independently decides the conversion rate $k$ based on the anticipated collection prices, while both participants engage in price competition for collection. This model serves as a baseline benchmark for subsequent comparative analyses with the biform game model.

\subsection{Two-stage biform game model}
\label{sec:biform_model}
In this subsection, we extend the non-cooperative game module to a non-cooperative and cooperative two-stage biform game model and explore its impact on economic benefits. Building on the scenario in the previous section, we assume the concrete manufacturer is willing to establish a R\&D investment partnership with the recycler. Based on this, we develop a waste concrete recycling management framework based on the biform game, combining non-cooperative and cooperative games to maximize economic benefits. Both participants agree to adopt this biform game structure to formulate their co-opetition relationships.

The biform game model we proposed for waste concrete recycling includes two stages (as illustrated in Fig. 2).

\textbf{The non-cooperative stage}: The manufacturer and recycler determine their pricing strategies via Nash Equilibrium. Notably, when making pricing decisions in this biform game, both participants anticipate how the final profits will be allocated through the cooperative game based on their R\&D investment cooperation. Therefore, this pricing strategy stage does not yield the final standalone profits but rather establishes the competitive environment.

\textbf{The cooperative stage}: The two participants engage in R\&D investment cooperation, where the manufacturer helps the recycler share a certain proportion of the R\&D investment. They can choose whether to form a cooperative alliance. To determine the profit distribution mechanism, we first calculate the characteristic values of the coalitions. Then, we use the Shapley value to allocate the cooperative surplus based on these characteristic functions. This analysis is conducted under the given competitive environment (i.e., the pricing strategy combination $(p_M, p_R)$) established in the non-cooperative stage. The profits allocated in the cooperative game will be treated as the objective functions for both participants in the non-cooperative stage, and we derive the equilibrium solutions by maximizing these profit functions.

\subsubsection{Cooperative game}
\label{sec:cooperative_game}

In this subsection, given the pricing strategy profile $(p_M, p_R)$ determined in the non-cooperative stage, we compute the characteristic value $V_{p_M,p_R}(\mathbb{C})$ for any possible coalition $\mathbb{C} \in \{\emptyset, \{M\}, \{R\}, \{M,R\}\}$ using the maximin principle. According to the maximin principle, a coalition $\mathbb{C}$ determines its own optimal decision for maximizing its profit, which is minimized by its opposing coalition. Thus, the characteristic function represents the minimum guaranteed profit a coalition can achieve when facing the worst-case decisions made by the opposing coalition.

For an empty coalition, the profit is certainly zero, thus $V_{p_M,p_R}(\emptyset) = 0$. Next, we calculate $V_{p_M,p_R}(\mathbb{C})$ for $\mathbb{C} \neq \emptyset$. Taking coalition $\{M\}$ as an example, since it acts as a single-player coalition, recycler $R$ forms the opposing coalition. The characteristic value $V_{p_M,p_R}(\{M\})$ depends on $M$'s cost-sharing decision $n$ and $R$'s conversion rate decision $k$, assuming the competitive environment $(p_M, p_R)$ is given. Thus, we apply the maximin principle as follows:
\begin{equation}
V_{p_M,p_R}(\{M\}) = \max_{n \in [0,1]} \min_{k \in [0,1]} \{\Pi_M(p_M, p_R, k, n)\}
\end{equation}
By solving this (detailed in Appendix B), we can obtain the characteristic value $V_{p_M,p_R}(\{M\})$. For other coalitions, we follow the same logic. We present $V_{p_M,p_R}(\mathbb{C})$ for all coalitions in Proposition 2.

\begin{prop}
\label{prop:characteristic}
Under any competitive situation $(p_M, p_R)$, the characteristic values of all possible coalitions $\mathbb{C}$ ($\mathbb{C} \neq \emptyset$) are given as follows:
\begin{equation}
V_{p_M,p_R}(\{M\}) = (s - p_M - w - \mu_M)(\alpha p_M - \beta p_R)
\end{equation}
\begin{equation}
\begin{aligned}
V_{p_M,p_R}(\{R\}) =& (s - p_R - m - \mu_R)(\alpha p_R - \beta p_M) + (w - m)(\alpha p_M - \beta p_R) \\
&+ \frac{\theta^2(\alpha p_R - \beta p_M)^2}{2c} - p_c[e(\alpha - \beta)(p_M + p_R) - G]
\end{aligned}
\end{equation}
\begin{equation}
\begin{aligned}
V_{p_M,p_R}(\{M,R\}) =& (s - p_M - m - \mu_M)(\alpha p_M - \beta p_R) + (s - p_R - m - \mu_R)(\alpha p_R - \beta p_M) \\
&+ \frac{\theta^2(\alpha - \beta)^2(p_M + p_R)^2}{2c} - p_c[e(\alpha - \beta)(p_M + p_R) - G]
\end{aligned}
\end{equation}
\end{prop}

 The proof is provided in Appendix B.

The coalitions' characteristic values computed in Proposition 2 constitute the cooperative game for R\&D investment cooperation. In cooperative game theory, the characteristic value reflects a coalition's bargaining power. To apply the Shapley value to determine a fair profit allocation scheme, it is crucial to check whether this cooperative game is convex and super-additive. If so, the core of this game is non-empty, and the Shapley value, as a single solution, will fall into the core.

\begin{prop}
\label{prop:convexity}
The cooperative game defined by the characteristic functions in Proposition 2 is convex and super-additive. Specifically, for any $(p_M, p_R)$, we have:
\begin{equation}
V_{p_M,p_R}(\{M,R\}) - V_{p_M,p_R}(\{M\}) - V_{p_M,p_R}(\{R\}) = \frac{\theta^2}{2c}(\alpha p_M - \beta p_R)[(\alpha p_M - \beta p_R) + 2(\alpha p_R - \beta p_M)] \geq 0
\end{equation}
Consequently, $V_{p_M,p_R}(\{M,R\}) \geq V_{p_M,p_R}(\{M\}) + V_{p_M,p_R}(\{R\})$, and the Shapley value allocation belongs to the core.
\end{prop}

 The proof is provided in Appendix C.

Proposition 3 implies that forming the grand coalition generates more benefits, providing a strong incentive for both the concrete manufacturer and the recycler to cooperate. 
The Shapley value is an important and widely accepted allocation mechanism in the cooperative game theory. Although cooperative games with a nonempty core offer several distribution strategies, this paper utilizes the Shapley value due to its superior practical acceptance among participants. This allocation method is effective because it aligns individual rewards with each player's distinct contribution to the overall coalition. The Shapley value allocations $\varphi_M$ and $\varphi_R$ are calculated as:

\begin{equation}
\varphi_M(p_M, p_R) = (s - p_M - w - \mu_M)(\alpha p_M - \beta p_R) + \frac{\theta^2(\alpha p_M - \beta p_R)[(\alpha p_M - \beta p_R) + 2(\alpha p_R - \beta p_M)]}{4c}
\end{equation}
\begin{equation}
\begin{aligned}
\varphi_R(p_M, p_R) =& (s - p_R - m - \mu_R)(\alpha p_R - \beta p_M) + (w - m)(\alpha p_M - \beta p_R) \\
&+ \frac{\theta^2[(\alpha p_R - \beta p_M)^2 + (\alpha p_M - \beta p_R + \alpha p_R - \beta p_M)^2]}{4c} \\
&- p_c[e(\alpha - \beta)(p_M + p_R) - G]
\end{aligned}
\end{equation}

In this sense, the players can achieve optimal profits in the cooperative game part, making it feasible to deploy the non-cooperative game by using these resulting allocations as payoff functions.

\subsubsection{Non-cooperative game}
\label{sec:noncooperative_game}

Following the biform game logic, the resulting allocations derived in the cooperative game part give the payoff functions in the non-cooperative game part. With the competitive environment defined by the pricing strategies, players anticipate how cooperation benefits will be allocated. Therefore, the payoff functions for $M$ and $R$ in the non-cooperative price competition stage are given by:
\begin{equation}
f_M(p_M, p_R) = \varphi_M(p_M, p_R)
\end{equation}
\begin{equation}
f_R(p_M, p_R) = \varphi_R(p_M, p_R)
\end{equation}

By using the Nash equilibrium concept, we maximize these two payoff functions. We take the first-order partial derivatives of $f_M(p_M, p_R)$ and $f_R(p_M, p_R)$ with respect to $p_M$ and $p_R$, respectively, and set them equal to zero. Solving these reaction functions (detailed in Appendix D) yields the equilibrium solutions as shown in Proposition 4.

\begin{prop}
\label{prop:biform}
Under the two-stage biform game model, the equilibrium collection prices $(p_M^{C*}, p_R^{C*})$ satisfy:
\begin{equation}
p_M^{C*} = \frac{
\displaystyle
\begin{aligned}
&4c^2 \left[ 2\alpha^2(s - w - \mu_M) + \alpha\beta(s - m - \mu_R) - \beta^2(w - m) \right] \\
&- 2c\theta^2 \Big\{ (\alpha-\beta) \left[ \alpha^2(s - w - \mu_M) - (\alpha^2 - \alpha\beta + \beta^2)(w - m) \right] - \alpha(\alpha^2 - \alpha\beta + \beta^2)(\mu_M - \mu_R) \Big\} \\
&- 4c\alpha p_c e (\alpha - \beta)(\alpha^2 - \alpha\beta + \beta^2)
\end{aligned}
}{
\displaystyle (\alpha-\beta)^3(\alpha+\beta)\theta^4 - 4c(\alpha-\beta)(3\alpha^2-\beta^2)\theta^2 + 4c^2(4\alpha^2-\beta^2)
}
\end{equation}
\begin{equation}
p_R^{C*} = \frac{
\displaystyle
\begin{aligned}
&4c^2\alpha \left[ 2\alpha(s - m - \mu_R) + \beta(s - 3w + 2m - \mu_M) \right] \\
&- 2c\alpha\theta^2 \Big\{ (\alpha-\beta) \left[ \alpha(w - m) + \beta(s - 3w + 2m - \mu_M) \right] + \alpha(\alpha - 2\beta)(\mu_M - \mu_R) \Big\} \\
&- 4c p_c e (\alpha - \beta) \left[ c(2\alpha^2 - \beta^2) - \theta^2\alpha(\alpha - \beta) \right]
\end{aligned}
}{
\displaystyle (\alpha-\beta)^3(\alpha+\beta)\theta^4 - 4c(\alpha-\beta)(3\alpha^2-\beta^2)\theta^2 + 4c^2(4\alpha^2-\beta^2)
}
\end{equation}
Substituting $(p_M^{C*}, p_R^{C*})$ back into the grand coalition's optimal decisions, the optimal conversion rate $k^{C*}$ and R\&D investment sharing proportion $n^{C*}$ are:
\begin{equation}
k^{C*} = \frac{\theta(\alpha-\beta)(p_M^{C*} + p_R^{C*})}{c}
\end{equation}
\begin{equation}
n^{C*} = \frac{2(\alpha p_R^{C*} - \beta p_M^{C*})^2 + 2(\alpha p_M^{C*} - \beta p_R^{C*})(\alpha p_R^{C*} - \beta p_M^{C*}) - (\alpha p_M^{C*} - \beta p_R^{C*})^2}{2 \big[ (\alpha-\beta)(p_M^{C*} + p_R^{C*}) \big]^2}
\end{equation}
\end{prop}

 The proof is provided in Appendix D.

By employing the Shapley value, we find a fair allocation solution which suggests reasonable profit functions for the non-cooperative game. The biform game enables us to derive analytical results for a complex problem where collection pricing competition and R\&D investment cooperation interact—a dynamic that is almost impossible to solve comprehensively using traditional single-stage game approaches.

\section{Equilibrium analysis}
\label{sec:equilibrium_analysis}
In this section, we conduct a comparative analysis of the equilibrium solutions between the non-cooperative model and the biform game model. By analyzing the optimal conversion rates, collection prices, market shares, and economic benefits across both scenarios, we reveal the theoretical mechanisms through which the biform game drives supply chain coordination. The outcomes of this analysis are articulated in Propositions 5 through 8.


\begin{prop}
\label{prop:technology_comparison}
Comparing the optimal RCA conversion rates between the non-cooperative game and the biform game, we have:
\begin{enumerate}[label={(\arabic*)}]
\item $k^{C*} > k^{N*}$;
\item $\frac{\partial k^{C*}}{\partial \theta} > \frac{\partial k^{N*}}{\partial \theta} > 0$; and $\frac{\partial k^{i*}}{\partial c} < 0$ for $i \in \{N, C\}$.
\end{enumerate}
\end{prop}

Proposition 5(1) indicates that the recycler will achieve a higher conversion rate under the biform game model. In the non-cooperative game model, the recycler bears the entire R\&D investment, leading to insufficient investment incentives due to a high marginal cost burden. The biform game introduces an R\&D investment sharing ($n^{C*}$), which internalizes the positive externalities of technological upgrades. By sharing the upfront costs, the manufacturer subsidizes the recycler's R\&D, incentivizing the recycler to extract more value from the waste concrete.

Proposition 5(2) demonstrates that while higher potential recycled value ($\theta$) stimulates technological upgrades in both models, the biform game exhibits greater sensitivity. When the recycled concrete becomes highly profitable the R\&D investment from the manufacturer helps the two enterprises upgrade the technology quickly. Without this financial help the recycler cannot afford expensive equipment upgrades even when the market is very lucrative.


\begin{prop}
\label{prop:pricing_comparison}
Regarding the collection pricing and resulting market shares, the following relationships hold:
\begin{enumerate}[label={(\arabic*)}]
\item Total collection quantity: $q_{total}^{C*} > q_{total}^{N*}$;
\item Manufacturer's market share advantage: $q_{M}^{C*} > q_{R}^{C*}$ and $q_{M}^{N*} > q_{R}^{N*}$ always hold, given $\mu_{M} < \mu_{R}$.
\end{enumerate}
\end{prop}

Proposition 6 explains the differences in the collection quantities of waste concrete. The first part shows that the total collection quantities is larger in the biform game. In the models without cooperation, the two enterprises simply fight over prices which reduces their profits and limits the total collection quantities. When they share the R\&D investment the manufacturer can benefit from the upgraded processing equipment of the recycler. This encourages the manufacturer to adjust its pricing strategy to attract more total waste concrete into the market rather than just fighting the recycler for the existing supply.

The second part of Proposition 6 proves that the concrete manufacturer always collects more waste than the recycler in both models. The manufacturer uses returning delivery trucks to transport waste which makes its transportation costs very low. The recycler must hire dedicated trucks which makes its transportation very expensive. Because the manufacturer spends much less on transportation it has more financial freedom to offer better collection prices to the market. This cost advantage allows the manufacturer to consistently capture a larger market share than the recycler.


\begin{prop}
\label{prop:profit_comparison}
Comparing the optimal profits of both participants across the two models, we have:
\begin{enumerate}[label={(\arabic*)}]
\item $\Pi_{M}^{C*} > \Pi_{M}^{N*}$ and $\Pi_{R}^{C*} > \Pi_{R}^{N*}$;
\item $\Pi_{M}^{C*} + \Pi_{R}^{C*} > \Pi_{M}^{N*} + \Pi_{R}^{N*}$.
\end{enumerate}
\end{prop}

Proposition 7 proves that establishing a cooperative alliance increases the financial returns for both enterprises. Proposition 7 (1) shows that the concrete manufacturer and the recycler both achieve higher individual profits in the biform game. Proposition 7 (2) confirms that the total profit of the entire recycling process is also larger. Although the manufacturer pays a portion of the R\&D investment. This shared effort improves the processing technology and expands the overall market scale. The expanded market generates enough additional revenue to benefit both participants. In summary, cooperation in the competitive environment provides higher profits for everyone compared to operating alone. Enterprises should actively seek cooperation.


\begin{prop}
\label{prop:carbon_impact}
Under the cap-and-trade mechanism, the impacts of the carbon trading price ($p_{c}$) on the equilibrium strategies are as follows:
\begin{enumerate}[label={(\arabic*)}]
\item $\frac{\partial p_{R}^{N*}}{\partial p_{c}} < 0$ and $\frac{\partial p_{R}^{C*}}{\partial p_{c}} < 0$;
\item $|\frac{\partial \Pi_{R}^{C*}}{\partial p_{c}}| < |\frac{\partial \Pi_{R}^{N*}}{\partial p_{c}}|$ (assuming total emissions exceed the free quota $G$).
\end{enumerate}
\end{prop}

Proposition 8 explains how the carbon trading price affects the recycling process. Proposition 8 (1) shows that a higher carbon price forces the recycler to lower its collection price in both models. When the price of carbon increases the recycler must pay more to process the waste. To manage this extra expense the recycler reduces the amount it pays to collect waste concrete.

Proposition 8 (2) reveals that cooperative investment enhances the recycler's resilience to carbon market volatility. While escalating carbon costs inherently erode profitability across both models, the cooperative game substantially decelerates this downward trend. Driven by shared R\&D investments, the improved processing technology yields higher-quality recycled aggregates, generating additional economic rents. This technological dividend empowers the recycler to better absorb the financial shocks of carbon compliance, proving that cooperation is vital for navigating the economic pressures of environmental regulations.

\section{Numerical simulation}
To illustrate the equilibrium outcomes more intuitively and investigate the impacts of key parameters (including price sensitivity coefficient, competition intensity, recycling value, and transportation costs) on collection prices, collection quantities, cost-sharing ratios, and technology improvement levels, we conduct numerical simulations using MATLAB 2023b.

\subsection{Parameter settings}
Waste concrete recycling represents a global challenge, and parameter values vary across different countries and contexts. To ensure the numerical simulation aligns with practical scenarios, we situate our analysis in the Chinese context. This selection is justified by two considerations: First, China has demonstrated strong policy commitment to waste concrete recycling, exemplified by initiatives such as the ``14th Five-Year Plan'' for solid waste pollution prevention and control, which explicitly mandates construction waste resource utilization rates to reach 60\% in large and medium-sized cities by 2025. Second, China's waste concrete recycling industry is at an early stage of development, with firms actively pursuing technological upgrades and seeking R\&D collaborations, which parallels the non-cooperative and cooperative game framework in this study. Accordingly, this context provides valuable insights for firm decision-making and government policy guidance.

Based on China's official policies, industry reports, project disclosures, expert interviews, and relevant literature, we calibrate the parameters as presented in Table~\ref{tab:parameters}.

\begin{table}[!htbp]
\centering
\captionsetup{font=normalsize, labelsep=period}
\caption{Parameter settings and baseline values}
\label{tab:parameters}
\normalsize
\begin{tabular}{lp{6cm}c}
\toprule
\textbf{Parameter} & \textbf{Source} & \textbf{Value} \\
\midrule
$\alpha$ & Drawing upon the Survey Report on the Resource Utilization of Construction Waste in Shanghai, the collection price for waste concrete typically ranges from 10 to 15 yuan per ton. Following the parameter specification in \citep{wang2025Mechanism}'s approach, we set $\alpha$ = 15/10 = 1.5 & 1.5 \\
\midrule
$\beta$ & According to the China Construction Waste Recycling Industry Overview, the compound annual growth rate of construction waste is approximately 13.4\%. Given the substantial capital requirements and technical barriers, the number of major firms remains relatively stable. We interpret the annual incremental waste concrete resources as the competition intensity. & 0.134 \\
\midrule
$m$ & Based on the report in World Environment journal, processing 100,000 cubic meters of waste concrete costs approximately 8 million yuan. We convert the unit. & 32 yuan/m$^3$ \\
\midrule
$w$ & Drawing from annual reports of A-share listed firms specializing in waste concrete recycling (e.g., Sinoma Resources, China Tianying, Shanghai Environment), the gross profit margins range from 7.88\% to 13.04\%. Taking the average margin of 9.84\%, we calculate $w = m \times (1 + 9.84\%) = 35$ yuan/m$^3$. & 35 yuan/m$^3$ \\
\midrule
$s$ & According to the 2025 project disclosure for the Fuzhou (Jiangxi Province) construction waste treatment franchise, the government provides a subsidy of 35 yuan per ton based on collection volume. We convert the unit. & 87.5 yuan/m$^3$ \\
\midrule
$\mu_M$ & We conducted market research by requesting quotes from multiple logistics providers and interviewing construction cost professionals. & 25 yuan/m$^3$ \\
\midrule
$\mu_R$ & We conducted market research by requesting quotes from multiple logistics providers and interviewing construction cost professionals. & 44 yuan/m$^3$ \\
\midrule
$c$ & Based on equipment pricing data from the China Construction Waste Recycling Industry Overview (including crushers, vibrating feeders, and magnetic separators), we calculated the weighted average cost coefficient. & 135,000 \\
\midrule
$\theta$ & According to Liu's survey report, the selling price of recycled concrete aggregate from waste concrete ranges from 50 to 60 yuan per ton. Taking the median value and converting units. & 137.5 yuan/m$^3$ \\
\midrule
$p_c$ & Referring to the carbon trading prices in various Chinese pilot markets in 2021, which ranged from 50 to 80 yuan per ton. Taking the median value. & 65 yuan/ton \\
\midrule
$G$ & We follow the parameter specification in the research of \cite{sun2021Optimal}. & 300 ton \\
\midrule
$e$ & According to the European Environment Agency (EEA), processing 1 ton of waste concrete generates approximately 0.02 to 0.05 tons of carbon dioxide emissions. Taking the median value. & 0.035 ton CO$_2$/m$^3$ \\
\bottomrule
\end{tabular}
\end{table}
\vspace{-15pt}

\subsection{Simulation results and analysis}
\label{sec:simulation_analysis}

To investigate the impacts of key parameters on equilibrium outcomes, we conduct a comprehensive numerical simulation analysis. Based on the baseline parameter values specified in Section 5.1, we vary each parameter within a reasonable range while holding other parameters constant. The simulation results are presented in Figs. 3-14.

\subsubsection{Comparison between Non-cooperative and Cooperative Games}
\label{sec:comparison}

This subsection presents the numerical comparison between the non-cooperative game model and the biform game model under varying key parameters. Our analysis focuses on how the cooperation affects the profits of both the manufacturer and the recycler under different market conditions.

Figure 3 illustrates the effects of the price sensitivity coefficient and the competition intensity on the optimal profits of the concrete manufacturer and the recycler. As the market becomes increasingly sensitive to price changes, the profits for both participants exhibit an upward trajectory across the non-cooperative and biform game models. Although high price sensitivity typically intensifies pricing battles and suppresses equilibrium prices, it simultaneously triggers a substantial expansion in total market demand. This increase in total collection quantities more than compensates for the reduced unit margins and drives up the total economic returns. Regardless of the fluctuations in price sensitivity, the cooperative strategy consistently achieves a Pareto improvement where both participants realize higher individual profits compared to the non-cooperative game model. The profit gap between the biform and non-cooperative game models widens as price sensitivity increases. In highly price-sensitive environments, the R\&D investment cooperation under the biform game lowers the marginal processing costs, which allows the allied enterprises to maintain robust profit margins. Furthermore, the profit allocation mechanism ensures that the recycler captures a larger share of the cooperative surplus. Because the recycler is the entity physically processing the waste to realize its high value, the Shapley value allocation rewards this specialized processing capacity and provides the recycler with a larger share of the expanded market profit.

Regarding the impact of competition intensity, the escalation of market competition exerts a steady downward pressure on the profits of both enterprises. This trend encapsulates the classic prisoner's dilemma within the front-end collection market, where fierce competition inevitably shrinks profit margins and forces participants to lower their collection prices. However, the R\&D investment cooperation demonstrates a robust buffering effect against these competitive frictions. Driven by shared economic objectives, the cooperative alliance prevents destructive price wars and allows both enterprises to sustain a higher profit baseline than they would operating independently. The surplus generated by cooperation remains relatively stable even under intense competition. The analysis further reveals a structural asymmetry in vulnerability between the two participants. The profit of the recycler demonstrates a heightened sensitivity to escalating competition relative to its operational scale. Without the protective umbrella of the cooperative alliance, the recycler is distinctly more vulnerable to external market shocks due to its specialized role and heavier technical cost burdens. Establishing a strategic alliance is therefore critical for the recycler to mitigate competitive risks and ensure long-term economic viability.

The analysis demonstrates that the biform game model captures the demand expansion caused by price sensitivity, which helps avoid the prisoner's dilemma in collection markets \citep{lin2025Cooperation}. The Shapley value allocation links the surplus distribution directly to each firm's operational contribution. This illustrates why the recycler's specialized processing capability is the primary driver for capturing market value, extending the fairness allocation principles of \cite{jia2023Novel}. Furthermore, the model demonstrates that forming a cooperative alliance prevents destructive price wars. By sharing R\&D investment, the recycler maintains stable profit margins even when front-end collection competition intensifies. This cooperative mechanism allows the recycler to survive severe competitive pressures while remaining an independent processing enterprise.

Figure 4 illustrates the influence of the the government subsidy and outsourcing fee on the economic outcomes of the two enterprises. As the subsidy level escalates, the overall financial returns for both supply chain participants exhibit a substantial upward trajectory. This indicates that targeted government financial interventions effectively expand the profit margins within the waste concrete recycling market. Counterintuitively, for concrete manufacturers, more subsidies are not necessarily better. Under low to moderate subsidy environments, the manufacturer derives greater profit from the cooperative alliance. Conversely, when the government subsidy surpasses a specific critical point, the non-cooperative strategy paradoxically becomes more lucrative for the manufacturer. In these highly subsidized scenarios, the financial obligation of sharing R\&D investment outweighs the cooperative surplus, effectively transforming the partnership into an economic burden for the manufacturer. In contrast, the recycler consistently realizes higher profits through the biform game across all subsidy scenarios. Furthermore, the cooperative premium for the recycler expands significantly as the subsidy increases, which reinforces the strong incentive for the specialized processor to maintain the strategic alliance when external policy support is robust.

The outsourcing fee represents a direct transfer payment from the manufacturer to the recycler, which results in opposed profit trajectories for the participants. As the outsourcing fee increases, the profit of the manufacturer experiences a continuous and severe decline due to the escalating operational costs associated with delegating the processing tasks. Similar to the subsidy dynamics, the manufacturer faces a inflection point regarding cooperation. When the outsourcing fee remains relatively low, the manufacturer benefits from the R\&D investment cooperation. However, as the transfer payment reaches elevated levels, the combined financial pressure of high service fees and shared technology investment costs renders the cooperative alliance less profitable than the non-cooperative game. Conversely, the profit of the recycler surges as the outsourcing fee increases. Throughout this transition, the recycler consistently secures a higher profit under the cooperative strategy compared to the non-cooperative baseline. This dynamic demonstrates that while high transfer payments secure the financial viability of the recycler, an excessive outsourcing fee destabilizes the cooperative willingness of the manufacturer and threatens the long-term sustainability of the joint innovation alliance.

Identifying a critical subsidy threshold challenges the assumption that higher subsidies always benefit recycling networks \citep{yuan2023Differentiated}. By modeling pricing competition and R\&D investment cooperation together, this analysis defines the specific boundary conditions for cooperative alliances. When external subsidies cross a certain threshold, the manufacturer's financial burden of sharing technology investments exceeds the cooperative surplus, making the alliance unprofitable for the manufacturer. Additionally, the analysis reveals a distinct inflection point for the outsourcing fee. Excessive transfer payments from the manufacturer to the recycler disrupt the balance between front-end collection and back-end treatment. This mechanism explains the instability of joint innovation alliances observed in third-party remanufacturing \citep{chen2025Exploring}, indicating that both policy interventions and inter-firm contracts require careful calibration.

Figure 5 illustrates the impact of the transportation cost differential and the potential recycled market value on the economic performance of both enterprises. As the transportation cost gap expands, the financial outcomes for the two participants diverge significantly. The profit of the manufacturer experiences a continuous increase because its inherent reverse logistics advantage becomes more pronounced in a high-friction logistical environment. By leveraging its established delivery network, the manufacturer dominates the front-end collection market and translates this logistical superiority into substantial economic returns. Consequently, the profit of the manufacturer under the cooperative alliance achieves the most significant growth. Conversely, the profit of the recycler suffers a steady decline as the escalating dedicated transportation costs compress its operational margins. When the logistical cost disadvantage is minimal, the recycler might prefer the non-cooperative strategy. However, as the transportation cost gap widens, the cooperative strategy yields higher profits for the recycler compared to operating independently. The specialized recycler becomes increasingly dependent on the cooperative alliance to buffer against severe logistical penalties. This proves that cooperation operates as an essential survival mechanism for the recycler in highly asymmetric logistical environments.

The analysis further examines the influence of the potential market value of the recycled concrete aggregate. An escalation in the market value elevates the profitability of both the concrete manufacturer and the recycler across both the non-cooperative and biform game models. A higher end-product value expands the total economic surplus available within the recycling system. Throughout this market expansion, the two participants in the biform game model consistently outperform them in the non-cooperative structure by enabling both participants to capture a larger portion of the augmented market wealth. Notably, when the market value is high, the profit growth of the recycler under the cooperative alliance is larger than in the non-cooperative game. Driven by the shared R\&D investments, the upgraded processing technology allows the recycler to capture the surging market value more efficiently. This direct link between high market value and increased cooperative willingness provides a clear direction for external interventions. If regulatory bodies implement measures to elevate the market acceptance and selling price of recycled products, such as green public procurement protocols or standardized certification systems, these actions will not only stimulate basic waste collection but also naturally induce enterprises to form voluntary research and development alliances to maximize their returns.

By incorporating the construction industry's transportation cost differences into the payoff functions, the simulation shows how cost advantages shape the collection market. A wider transportation cost gap forces the recycler to rely on the cooperative alliance to maintain profitability. This links facility integration strategies \citep{liu2019Expand} with coopetition theory \citep{lin2025Cooperation}. Furthermore, while previous literature notes a positive relationship between end-product valuation and recycling viability \citep{li2020Key}, the biform game approach specifies the mechanism. It demonstrates that higher market values directly motivate voluntary R\&D partnerships. Consequently, policy-driven market enhancements \citep{su2024Stakeholder}., such as green procurement standards, can shift the burden of technological upgrades from government mandates to shared market investments.


\subsubsection{Impact of Cap-and-Trade Mechanism}
\label{sec:cap_and_trade}

This section evaluates how the cap-and-trade mechanism impacts the decisions and economic benefits of the waste manufacture and the recycler within the waste concrete recycling process. We focus on three parameters of the policy, which are the carbon trading price, the carbon emission coefficient, and the government allocated carbon quota.

Figures 6(a) and 6(b) illustrate the impacts of the carbon trading price and the actual carbon emissions per unit on the optimal profits of the concrete manufacturer and the recycler, respectively. As both the carbon trading price and the unit carbon emissions escalate, the economic returns for both participants exhibit a downward trajectory. Notably, the profit of the recycler is more sensitive to the escalating carbon trading price, experiencing a much steeper and more profound decline compared to the manufacturer. This reflects the recycler's direct operational exposure and heavy financial burden regarding emission compliance. While an increase in unit carbon emissions also erodes profits across the board, its overall negative impact is milder than that of the carbon price, though it still affects the specialized recycler more distinctly than the manufacturer.

The comparative analysis of the two game reveals the advantage of the cooperative alliance. Under the biform game model, both the magnitude and the rate of profit decline for the two enterprises are smaller than those observed in the non-cooperative game. This reveals that R\&D investment cooperation effectively acts as a robust financial buffer. By establishing a cooperative alliance, the concrete manufacturer and the recycler can mitigate the negative financial shocks induced by volatile carbon prices and stringent emission penalties, thereby preserving their profit margins and enhancing the economic resilience of the entire waste concrete recycling network.

Fig. 7 demonstrates how the participants reacts when the carbon trading price increases. As the price goes up, the recycler faces a much higher cost to process the waste concrete because they are the actual party generating emissions. To balance this extra expense, the recycler has to reduce the collection price offered to the market. The concrete manufacturer does not bear the emission cost, so their collection price stays stable. The collection quantities also drop and it pull down the total quantities. Furthermore, the heavier financial burden leaves the recycler with less capital to share the R\&D investment. This leads to a lower conversion rate. Ultimately, the rising carbon price cuts into the profit margins of both enterprises, but the recycler suffers a much sharper decline due to their direct exposure to the carbon market.

These findings align with prior research showing that carbon pricing mechanisms alter the cost structure of the supply chains. Under cap-and-trade mechanism, enterprises directly exposed to emission compliance costs bear disproportionate burdens, creating structural inequalities in inter-firm cooperation \citep{bai2022Distributionally}. The differential vulnerability between manufacturers and recyclers mirrors observations in EV battery closed-loop supply chains where carbon costs reshape pricing dynamics across recycling channels \citep{wu2024Electric}. Our results extend this insight by demonstrating that rising carbon prices compress margins and undermine the financial prerequisites for cooperative R\&D, reducing both the conversion rate and the R\&D investment sharing proportion. This suggests that in high carbon-price environments, policy interventions may be necessary to maintain cooperative incentives in recycling networks \citep{huang2024Optimal}.


Fig. 8 illustrates the impact of the carbon emission coefficient. A higher emission coefficient means the recycler produces more carbon dioxide for every unit of waste processed. This inflates the total carbon compliance cost in the same way a higher trading price does. Facing higher operational costs, the recycler scales back their front end procurement price. The shrinking market volume and increased cost pressure make both participants less willing to achieve R\&D investment cooperation, causing a steady decline in the conversion rate and the sharing ratio. The profit trajectories confirm that if the processing technology is not environmentally friendly, the resulting carbon costs will damage the economic benefits of cooperation in the waste recycling process.

This finding agrees with previous studies on emission reduction in waste management. When the processing technology generates a large amount of carbon emissions, the recycler faces high compliance costs. These high costs leave the recycler with less money for R\&D \citep{besklubova2026Establishing}. Because processing waste concrete becomes very expensive under high emission levels, both the concrete manufacturer and the recycler lose the economic incentive to invest in R\&D. Consequently, the conversion rate and the investment sharing proportion paid by the recycler both decrease. The results also show that the R\&D investment cooperation is easily broken by high carbon costs. The enterprises can easily lower their collection prices to survive sudden cost increases, but they will quickly stop their joint R\&D investments when profit margins shrink. This proves that developing processing technologies with low carbon emissions is strictly necessary to keep the cooperative partnership profitable and stable \citep{chen2025Lowcarbon}.

Fig. 9 focuses on the relationship between the government allocated carbon quota and the optimal profit of the recycler. The graph shows an upward trend. The free carbon quota acts as a financial buffer. When the quota is larger, the recycler is allowed to emit more carbon for free and needs to purchase fewer allowances from the trading market. This cost saving translates into the recycler profit.
However, this presents a policy dilemma. While excessive free allocation protects recycler profits, it may undermine the environmental objectives of the cap-and-trade mechanism. If quotas are too generous, the carbon price signal weakens, reducing incentives for recyclers to invest in low-carbon technologies. Indeed, research shows that over-allocation can create a "green paradox" where firms rely on free permits rather than pursuing emission reductions \citep{huang2024Optimal}. Therefore, the optimal quota design must balance maintaining economically viable recycling operations and preserving sufficient carbon pricing pressure to drive technological upgrades. The key is to set quotas at levels that prevent financial distress but still motivate R\&D investment \citep{peng2022If}.

In summary, the cap-and-trade mechanism introduces an environmental cost that reshapes the cost structure of the concrete manufacturer and the recycler. Because the recycler processes the waste concrete, they are the most vulnerable to high carbon prices and poor emission efficiencies. When these carbon costs become too heavy, the recycler is forced to shrink their collection quantities and reduce R\&D investments, which negatively impacts the entire waste concrete recycling supply chain. The free carbon quota can effectively relieve this pressure. Therefore, to ensure a healthy and sustainable recycling industry, it is crucial to find a balance between pushing for emission reductions and providing enough quota buffers to maintain the economic motivation of the recycling enterprises.

\section{Conclusion and implications}
\subsection{Conclusions}

This study develops a two-stage biform game model to investigate the coopetition relationships between the concrete manufacturer and the recycler in the waste concrete recycling process under the cap-and-trade mechanism. By comparing the equilibrium solutions of the biform game model with a non-cooperative game model, the study draws the following principal conclusions:

First, the results demonstrate that cooperation in the competition environment is more beneficial for both participants than the pure non-cooperative game model. Cooperation incentivizes a higher conversion rate and leads to an expansion of the total collection quantities. And this achieves a Pareto improvement, where both the manufacturer and the recycler realize higher individual profits compared to the non-cooperative game. Furthermore, the equilibrium analysis reveals that financial incentives are not simply "the more, the better." Counterintuitively, when government subsidies surpass a specific critical point, the non-cooperative strategy becomes more lucrative for the manufacturer. A similar inflection point exists for operational costs; excessive outsourcing fees destabilize the manufacturer's willingness to cooperate.

Second, concerning the impacts of the cap-and-trade mechanism, the study finds that an increase in the carbon trading price generally raises costs and reduces profit margins of the concrete manufacture and the recycler. Cooperation in the competition environment provides a buffering effect against such volatility. By sharing the R\&D investment, the manufacturer enables the recycler to reach a higher conversion rate, which generates greater marginal revenue and enhances the whole supply chain's resilience to rising carbon market costs. Furthermore, the government-allocated free carbon quota serves as an essential financial buffer to maintain the profitability of specialized processors. 

\subsection{Theoretical implications}
This study contributes to the related literature on construction and demolition waste management and low-carbon recycling supply chains in two ways:

First, this study constructs a two-stage biform game model for waste concrete recycling, expanding the application boundaries of biform game theory. Rather than treating the collection and processing stages as isolated activities, this model captures the intertwined dynamics of front-end pricing competition and back-end technology investment cooperation. This overcomes the limitations of traditional single-stage non-cooperative games. The model provides a concrete analytical paradigm for understanding how the waste manufacturer and the recycler can achieve cooperation within the competitive environment.

Second, this research innovatively integrates the carbon emission reduction and cap-and-trade mechanism into the economic decision of waste concrete recycling management. Previous operational models generally treat environmental policies merely as exogenous cost penalties. By coupling carbon market dynamics with cooperative investment strategies, our research theoretically demonstrates how strategic alliances can transform carbon compliance from a static constraint into a driver for R\&D investment. Furthermore, this study quantifies the buffering effect of cooperation against carbon price volatility, thereby advancing the theory of low-carbon supply chain resilience.

\subsection{Practical implications}

First, concrete manufacturers and specialized recyclers should establish cooperative alliances. The results demonstrate that this is the optimal path for maximizing both systemic and individual profits. To operationalize this alliance effectively, firms should transition from traditional transactional outsourcing to deep strategic integration. Specifically, partners should design KPI-linked cost-sharing contracts. Rather than a simple, flat capital injection, the manufacturer's financial support should be contractually tied to the recycler's technological performance metrics, such as the aggregate conversion rate. Furthermore, the alliance should implement a joint logistics dispatch system. By integrating data across the supply chain, the recycler can coordinate its processing schedule with the manufacturer's delivery routes, fully exploiting the manufacturer's reverse logistics advantage. This operational integration ensures a stable, large-scale supply of waste concrete, thereby securing the economies of scale necessary for back-end profitability.

Second, participants in the waste concrete recycling supply chain should pay close attention to carbon emissions and the cap-and-trade mechanism, treating them as key operational variables rather than mere compliance costs. To effectively respond to carbon market volatility, concrete manufacturers and recyclers are encouraged to explore an internal carbon risk-sharing and dividend mechanism. When negotiating the outsourcing fee and the technology investment ratio, partners can dynamically factor in the projected carbon trading prices. For instance, the cooperative contract could include a floating adjustment clause: during periods of high carbon prices, the manufacturer temporarily increases its cost-sharing proportion or transfer payments to buffer the recycler's compliance burden; conversely, any financial dividends gained from selling surplus carbon quotas—achieved through the jointly funded high-efficiency waste concrete processing technology—could be shared between both parties. Additionally, partners can consider implementing joint carbon accounting across the waste concrete collection and processing stages to accurately track and certify emission reductions, helping to optimize their returns in the external carbon market.

Third, government policymakers can shift from uniform industry incentives to targeted structural interventions. To address the inherent capability asymmetry in the market, authorities should implement specialized logistics subsidies exclusively for dedicated recyclers to neutralize their transportation cost disadvantages compared to manufacturers. Furthermore, policymakers must recognize the critical threshold effect of financial interventions. As our findings indicate that financial incentives are not simply "the more, the better," regulators must carefully calibrate subsidy levels. Subsidies should be dynamically evaluated and kept within an optimal range that stimulates market activity without disrupting the delicate financial balance required for strategic R\&D alliances. Additionally, regulators need to maintain a predictable carbon market environment with stable price guidance, which provides enterprises with the long-term confidence required to commit to capital-intensive cooperative investments.

\section{Limitations and future research directions}
There are certain limitations that highlight avenues for future research. First, this study primarily focuses on a basic dual-channel supply chain consisting of a single traditional manufacturer and a single specialized recycler. However, the actual waste concrete recycling ecosystem often involves more complex interactions among multiple stakeholders. Future research could expand the current modeling framework to incorporate multiple competing recyclers, third-party logistics providers, or explicitly consider the downstream demand market's acceptance of recycled materials, thereby more comprehensively revealing the operational mechanisms of a closed-loop supply chain network with multi-party participation.

Second, this research employs a static game framework to analyze the equilibrium decisions and profit distribution of supply chain members. Future research could integrate differential games or evolutionary game theory to explore the dynamic evolutionary trajectories of enterprises' technology R\&D investments over multiple periods, as well as the long-term stability conditions of the outsourcing cooperation.

\section*{Data availability statement}
All data, models, or codes that support the findings of this study are available from the corresponding author upon reasonable request.

\section*{Acknowledgments}
This paper was supported by the National Natural Science Foundation of China (Project Numbers: 72371189, 72371190) and the program of China Scholarship Council (CSC) (No. 202206260227).

\section*{Supplemental Material}
Appendix are available online at https://doi.org/10.5281/zenodo.20528901

\bibliographystyle{elsarticle-harv}
\bibliography{references}


\end{document}
